\section{Convergence and Performance Analysis}
\label{app:convergence}
\subsection*{Spatio-temporal convergence and performance}

We verify spatial consistency using field–norm error estimates on nested meshes with non-matching transfer. All quantities are dimensional; averages are normalized by $|\Omega|$.

\paragraph{Meshes and notation.}
Let $\{\mathcal{T}_i\}_{i=1}^L$ be a family of tetrahedral meshes on $\Omega$ obtained by uniform refinement, with characteristic sizes $h_i=1/N_i$ (we report $h_i=1/N_i$ explicitly). On each $\mathcal{T}_i$ we solve the fully coupled problem to a fixed time $T$ with \emph{the same} $\Delta t$ and model parameters; this avoids additional time–discretization bias from remounting left–hand sides. Denote by
\[
\mathbf{u}_i,\quad S_i,\quad \rho_i,\quad \mathbf{A}_i
\]
the discrete fields on level $i$ in their conforming $H^1$ spaces $V_{h_i},Q_{h_i},Q_{h_i},T_{h_i}$. Because the solution can exhibit sharp internal layers (bone/void fronts), we do not assume high regularity.

\paragraph{Field–norm pair errors on the finer mesh.}
For two consecutive levels $(H,h)=(i,i{+}1)$ we evaluate errors on the finer mesh in a \emph{higher–order observation space}. For any field $X\in\{\mathbf{u},S,\rho,\mathbf{A}\}$ we construct $W_{h}^{(X)}$ by raising the polynomial degree by $\Delta p=2$, transfer $X_H$ to $W_{h}^{(X)}$ by non-matching interpolation, and interpolate the fine-field values into the same space. With $\mathcal{I}_{H\to h}$ denoting the non-matching interpolation and $\widetilde X_h$ the interpolated fine solution, the pair error reads
\[
e_{H\to h}^{X,\;L^2} \;=\; \big\| \widetilde X_h - \mathcal{I}_{H\to h} X_H \big\|_{L^2(\Omega)},
\qquad
e_{H\to h}^{X,\;H^1\text{-semi}} \;=\; \big\| \nabla(\widetilde X_h - \mathcal{I}_{H\to h} X_H) \big\|_{L^2(\Omega)} .
\]
Vector/tensor seminorms use the canonical componentwise sum. All integrals are evaluated with sixth-order quadrature on $\mathcal{T}_h$. For a sequence of pairs we obtain $\{(h_{i+1},e_{i\to i+1}^{X,\bullet})\}_{i=1}^{L-1}$ and estimate an experimental order of convergence (EOC) $p$ from the log–log fit
\[
\log e \approx p\,\log h + \mathrm{const},\qquad
p \;=\; \arg\min_{p} \sum_{i} \big(\log e_{i\to i+1}^{X,\bullet} - p\log h_{i+1} - c\big)^2 ,
\]
reported for both $L^2$ and $H^1$–semi norms. When only three levels are available, we also quote the two–ratio estimate $p \approx \log(e_{H\to h}/e_{h\to h/2})/\log(r)$ with $r=h_H/h_h$.

\begin{figure}[htbp]
  \centering
  \IfFileExists{images/convergence_plot.png}{%
    \includegraphics[width=1\linewidth]{images/convergence_plot.png}%
  }{\fbox{\parbox{1\linewidth}{\centering Missing image: \texttt{images/convergence_plot.png}}}}
  \caption{Convergence of spatial and temporal errors under mesh and time-step refinement. Solid lines: errors (norms specified in the text). Dotted lines: ideal slopes for $P^1$ elements and the Backward--Euler scheme. Shaded bands show interquartile ranges across regions of interest.}
  \label{fig:individual_convergence}
\end{figure}
